\documentclass[conference]{IEEEtran}
\usepackage{cite}
\usepackage{amsmath,amssymb,amsfonts}
\usepackage{algorithmic}
\usepackage{graphicx}
\usepackage{textcomp}
\usepackage{xcolor}
\usepackage{hyperref}
\usepackage{booktabs}
\usepackage{listings}
\usepackage{algorithm}
\usepackage{url}

\lstset{
  basicstyle=\ttfamily\footnotesize,
  breaklines=true,
  frame=single,
  numbers=left,
  numberstyle=\tiny,
  tabsize=2
}

\def\BibTeX{{\rm B\kern-.05em{\sc i\kern-.025em b}\kern-.08em
    T\kern-.1667em\lower.7ex\hbox{E}\kern-.125emX}}

\begin{document}

\title{Cognitive Harness: A Plugin-Based LLM Reasoning Engine with Metacognitive Loop Recovery and Evolutionary Prompt Optimization}

\author{\IEEEauthorblockN{Saptarshi E.}
\IEEEauthorblockA{%
Independent Researcher}}

\maketitle

\begin{abstract}
Large Language Models (LLMs) frequently enter unproductive reasoning loops during multi-step problem solving, exhibiting hedging language and repeated phrases that waste computational resources without advancing toward a solution. We present \textbf{Cognitive Harness}, an open-source, plugin-based reasoning engine that addresses this limitation through four integrated mechanisms: (1)~a \emph{metacognitive stuck detector} that identifies stalled reasoning via heuristic analysis of hedging patterns and n-gram repetition; (2)~a \emph{parallel sub-agent spawner} that concurrently explores alternative solution strategies when the primary agent stalls; (3)~a \emph{two-tier context compaction} system inspired by Lean-style proof compaction that maintains reasoning coherence within bounded context windows; and (4)~an \emph{evolutionary prompt optimization} framework (EGPA) that uses genetic algorithms with LLM-driven mutation to discover higher-performing system prompts over successive generations. The system further integrates a hybrid memory subsystem combining ChromaDB vector embeddings with a knowledge graph, fused via Reciprocal Rank Fusion (RRF), and exposed through the Model Context Protocol (MCP). Cross-session global memory and a working-memory scratchpad enable persistent knowledge accumulation across interactions. A modular LangGraph-based orchestration layer allows all components to be independently replaced or extended. We describe the system architecture comprising over 4,000 lines of Python across 40+ modules, present its plugin registration mechanism, and discuss the design trade-offs underlying each subsystem. Evaluation on a 12-test end-to-end verification suite covering vault operations, graph queries, embedding search, stuck detection, and full LLM reasoning demonstrates correct functional behavior across all components.
\end{abstract}

\begin{IEEEkeywords}
large language models, reasoning engines, metacognition, prompt optimization, genetic algorithms, knowledge graphs, context management, multi-agent systems
\end{IEEEkeywords}

\section{Introduction}
\label{sec:intro}

The emergence of Large Language Models (LLMs) as general-purpose reasoning tools has created a fundamental tension: while models such as GPT-4~\cite{achiam2023gpt4}, DeepSeek-V3~\cite{deepseekai2024deepseekv3}, and Claude~\cite{anthropic2024claude} exhibit remarkable chain-of-thought capabilities, they remain prone to \emph{reasoning stalls}---sustained periods where the model generates plausible-sounding but non-progressive output. These stalls manifest as hedging language (``I'm not sure'', ``Let me try again''), repeated phrases across successive responses, and circular reasoning patterns that consume context window capacity without advancing toward a solution~\cite{wei2022chain}.

Prior work has addressed individual aspects of this problem. LangChain~\cite{langchain2023} and LangGraph~\cite{langgraph2024} provide orchestration primitives for multi-step LLM workflows. DSPy~\cite{khattab2023dspy} introduces bootstrapping-based prompt optimization. ReAct~\cite{yao2023react} interleaves reasoning and action. Reflexion~\cite{shinn2023reflexion} adds verbal self-reflection. However, no existing framework integrates \emph{metacognitive loop detection}, \emph{parallel recovery}, \emph{context compaction}, and \emph{evolutionary prompt improvement} into a single cohesive architecture that operates autonomously during multi-turn conversations.

This paper presents \textbf{Cognitive Harness}, a plugin-based LLM reasoning engine designed to maintain productive reasoning over extended interactions. Our key contributions are:

\begin{enumerate}
    \item A \textbf{metacognitive stuck detector} that identifies reasoning stalls through heuristic analysis of hedging language patterns and 4-gram repetition across the last five assistant messages, achieving detection without additional LLM calls.
    \item A \textbf{parallel sub-agent spawner} that, upon stuck detection, concurrently launches $N$ independent reasoning agents (default $N{=}3$), each assigned a distinct problem-solving strategy, to explore the solution space in parallel.
    \item A \textbf{two-tier context compaction} mechanism that uses LLM-assisted summarization to maintain reasoning coherence within bounded context windows, operating at both working-buffer and full-state granularities.
    \item An \textbf{Evolutionary Genetic Prompt Algorithm (EGPA)} that represents system prompts as genomes of discrete genes (role, constraints, style), evolves them through tournament selection, uniform crossover, and LLM-driven mutation, and selects winners via a composite fitness function.
    \item A \textbf{hybrid memory system} combining ChromaDB vector similarity search with a knowledge graph queried via extracted keywords, fused using Reciprocal Rank Fusion (RRF) with configurable $\alpha$ and $\beta$ weights, and served through a Model Context Protocol (MCP) server.
\end{enumerate}

The system is implemented in approximately 4,000 lines of Python across 40+ modules, uses LangGraph for orchestration, LiteLLM for model-agnostic LLM access, and is designed so that every component can be independently replaced via the plugin registry.

\section{Related Work}
\label{sec:related}

\textbf{LLM Orchestration Frameworks.} LangChain~\cite{langchain2023} provides chaining primitives for LLM applications, while LangGraph~\cite{langgraph2024} extends this with stateful, graph-based workflows supporting cycles and conditional routing. AutoGen~\cite{wu2023autogen} enables multi-agent conversations but does not address metacognitive monitoring. CrewAI~\cite{crewai2024} coordinates role-based agents but lacks adaptive recovery from reasoning stalls.

\textbf{Prompt Optimization.} DSPy~\cite{khattab2023dspy} optimizes LLM programs through bootstrapping and signature-based compilation. PromptBreeder~\cite{fernando2023promptbreeder} uses evolutionary strategies for prompt improvement. APE~\cite{zhou2023large} employs LLMs as prompt generators. Our EGPA approach shares the evolutionary spirit of PromptBreeder but introduces structured gene decomposition (role, constraints, style) and a composite fitness function incorporating success rate, latency, and token efficiency.

\textbf{Reasoning Improvement.} Chain-of-Thought prompting~\cite{wei2022chain} demonstrates emergent reasoning through few-shot exemplars. Tree-of-Thoughts~\cite{yao2024tree} explores branching reasoning paths. Reflexion~\cite{shinn2023reflexion} adds verbal self-reflection after failures. Self-Refine~\cite{madaan2023selfrefine} iteratively improves outputs through self-feedback. Cognitive Harness differs by detecting stuck states \emph{autonomously} rather than requiring explicit failure signals, and by recovering through parallel exploration rather than sequential reflection.

\textbf{Memory-Augmented LLMs.} MemGPT~\cite{packer2023memgpt} manages virtual context through memory pagination. Retrieval-Augmented Generation (RAG)~\cite{lewis2020rag} retrieves relevant documents before generation. Our hybrid retrieval approach combines dense (embedding) and sparse (knowledge graph) retrieval with RRF fusion~\cite{cormack2009reciprocal}, and persists knowledge across sessions through an Obsidian-style vault with YAML frontmatter.

\section{System Architecture}
\label{sec:architecture}

\subsection{Overview}

Cognitive Harness is structured as a layered architecture with five principal subsystems: the \emph{Orchestration Layer}, the \emph{Plugin Graph}, the \emph{Chat Engine}, the \emph{Memory Subsystem}, and the \emph{Self-Improvement Module}. Fig.~\ref{fig:architecture} illustrates the high-level data flow.

\begin{figure}[htbp]
\centerline{
\fbox{\parbox{0.9\columnwidth}{\footnotesize
\textbf{System Architecture}\\[4pt]
User $\rightarrow$ \textsc{Chat Engine} $\leftrightarrow$ \textsc{LLM Client} (LiteLLM)\\
$\downarrow$ tool calls\\
\textsc{Tool Registry} $\rightarrow$ Built-in Tools $|$ MCP Tools\\
$\downarrow$ reasoning loop\\
\textsc{LangGraph Orchestrator}:\\
\quad Thinker $\rightarrow$ Stuck Detector $\rightarrow$ \{Sub-Agent Spawner\\
\quad $\rightarrow$ Compactor\} $\rightarrow$ Context Manager $\rightarrow$ [loop$|$END]\\
$\downarrow$ memory\\
\textsc{Memory}: Vault $+$ KnowledgeGraph $+$ ChromaDB\\
$\downarrow$ optimization\\
\textsc{EGPA}: Tracker $\rightarrow$ Optimizer $\rightarrow$ Evolution Engine
}}}
\caption{Cognitive Harness system architecture showing the five principal subsystems and their interactions.}
\label{fig:architecture}
\end{figure}

\subsection{Orchestration Layer}

The orchestration layer uses LangGraph's \texttt{StateGraph} to define a directed graph of plugin nodes connected by conditional edges. The shared state is represented as a typed dictionary (\texttt{HarnessState}) containing:

\begin{itemize}
    \item \texttt{checkpoints}: A list of verified reasoning derivations, analogous to lemmas in a proof system.
    \item \texttt{working\_buffer}: The active message history between the user and the model.
    \item \texttt{is\_stuck}: A boolean flag set by the stuck detector.
    \item \texttt{sub\_agent\_results}: Outputs from parallel sub-agents.
    \item \texttt{iteration}/\texttt{max\_iterations}: Loop control for bounded execution.
\end{itemize}

The graph topology is:
\begin{equation}
\text{Thinker} \rightarrow \text{StuckDetector} \xrightarrow{\text{stuck?}} \begin{cases} \text{SubAgents} \rightarrow \text{Compactor} & \text{if stuck} \\ \text{ContextMgr} & \text{otherwise} \end{cases}
\end{equation}

The context manager conditionally routes back to the thinker or terminates based on iteration count and stuck status.

\subsection{Plugin System}

All reasoning components inherit from \texttt{BaseNode}, an abstract class requiring a single method:

\begin{lstlisting}[language=Python, caption={Plugin base class interface.}]
class BaseNode(ABC):
    name: str
    @abstractmethod
    async def process(
        self, state: HarnessState,
        llm: LLMClient
    ) -> HarnessState: ...
\end{lstlisting}

Plugins self-register via the \texttt{@register\_node("name")} decorator, which inserts the class into a global registry dictionary. The orchestrator discovers plugins at graph-build time by importing their modules, avoiding configuration files in favor of code-level declaration.

\subsection{Chat Engine}

The \texttt{ChatEngine} manages multi-turn conversations with a tool execution loop and optional streaming. For each user message, it:

\begin{enumerate}
    \item Constructs a dynamic system prompt incorporating: the base instruction, available skills, scratchpad entries, relevant global memory entries, and MCP memory context.
    \item Calls the LLM via LiteLLM with OpenAI-format tool definitions.
    \item Executes any returned tool calls through the \texttt{ToolRegistry}.
    \item Loops until the model returns a text-only response (no tool calls).
    \item Records the interaction metric in the prompt tracker.
\end{enumerate}

The system prompt is dynamically optimized: if the prompt optimizer identifies a higher-scoring prompt variant, it is automatically substituted.

\subsection{Tool System}

The \texttt{ToolRegistry} provides a unified interface for built-in tools (file operations, shell execution, Git operations, web search, skill management) and MCP-sourced tools. Each tool is registered with a name, description, JSON Schema parameters, and a handler function. The registry generates OpenAI function-calling format tool definitions and injects human-readable descriptions into the system prompt.

\section{Metacognitive Stuck Detection}
\label{sec:stuck}

The stuck detector analyzes the last $k$ assistant messages (default $k{=}5$) for two categories of stall indicators:

\textbf{Hedging Pattern Detection.} Seven regex patterns capture common hedging language:

\begin{lstlisting}[language=Python, caption={Hedging patterns for stuck detection.}]
HEDGING_PATTERNS = [
    r"i'?m not sure",
    r"maybe",
    r"let me try (again|something else)",
    r"i don'?t (know|understand)",
    r"this isn'?t working",
    r"i'?m stuck",
    r"going in circles",
]
\end{lstlisting}

A message is flagged as hedging if $\geq 2$ distinct patterns are detected in the concatenation of the last 3 messages. This threshold was empirically calibrated to balance sensitivity against false positives from legitimate uncertainty expressions.

\textbf{N-gram Repetition Analysis.} The detector extracts all 4-grams from each of the last 5 assistant messages (case-insensitive), counts their frequencies using a \texttt{Counter}, and flags as stuck if any 4-gram appears $\geq 3$ times across messages. This captures the ``going in circles'' phenomenon where the model rephrases the same idea repeatedly.

The stuck detector requires no additional LLM calls---it operates purely on the existing message buffer, adding negligible latency ($<$1ms per invocation).

\section{Parallel Sub-Agent Recovery}
\label{sec:subagent}

Upon stuck detection, the system spawns $N$ (default 3) parallel sub-agents via \texttt{asyncio.gather}. Each sub-agent receives:
\begin{itemize}
    \item The current checkpoint context (all verified derivations).
    \item The specific response where the main agent stalled.
    \item A \emph{distinct} exploration strategy from a predefined set (e.g., ``investigate caching'', ``check for unused parameters'', ``look for circular references'').
\end{itemize}

Sub-agents execute concurrently, and their results---including success/failure status and proposed solutions---are collected with exception isolation so that individual failures do not halt recovery. The \texttt{CompactorNode} then synthesizes successful sub-agent outputs into a new verified checkpoint using a structured Premise$\rightarrow$Step$\rightarrow$Conclusion format, discarding conversational filler and unverified claims.

This approach has a key advantage over sequential recovery strategies (e.g., Reflexion~\cite{shinn2023reflexion}): wall-clock recovery time is bounded by the slowest sub-agent rather than the sum of all attempts.

\section{Two-Tier Context Compaction}
\label{sec:compaction}

Context window management is critical for sustained reasoning. The \texttt{ContextManagerNode} implements a two-tier compaction strategy:

\textbf{Tier~1: Working Buffer Compaction} ($\tau_1 = 4000$ tokens). When the working buffer exceeds $\tau_1$, it is summarized into a new checkpoint via an LLM call with a compaction-specific prompt that enforces Premise$\rightarrow$Steps$\rightarrow$Conclusion format. The buffer is then replaced with a single ``Context compacted. Continue reasoning.'' message.

\textbf{Tier~2: Full Compaction} ($\tau_2 = 12000$ tokens). When total context (all checkpoints plus working buffer) exceeds $\tau_2$, all content is collapsed into a single, dense, self-contained derivation. All prior checkpoints are replaced with this unified derivation.

This two-tier design ensures that token usage grows logarithmically with reasoning depth rather than linearly, while the structured checkpoint format preserves logical chains.

\section{Evolutionary Prompt Optimization}
\label{sec:evolution}

The self-improvement subsystem consists of three integrated components.

\subsection{Prompt Tracker}

Every prompt--response pair is recorded as a \texttt{PromptMetric} containing: prompt ID, prompt text, response text, tokens used, latency, tool call count, and success status. A composite fitness score $S$ is computed as:
\begin{equation}
S = 0.6 \cdot R_s + 0.2 \cdot (1 - \hat{L}) + 0.2 \cdot (1 - \hat{T})
\label{eq:fitness}
\end{equation}
where $R_s$ is the success rate, $\hat{L} = \min(L_{avg} / 10000, 1)$ is the normalized average latency (ms), and $\hat{T} = \min(T_{avg} / 4000, 1)$ is the normalized average token usage. This multi-objective formulation balances correctness (60\%), responsiveness (20\%), and efficiency (20\%).

\subsection{Prompt Optimizer}

The optimizer monitors tracked metrics and initiates optimization cycles when sufficient samples accumulate ($\geq 10$ by default). It generates $n$ prompt variations by querying the LLM with performance context, then tracks each variation's real-world performance. Winners are selected by comparing composite scores, with a 5\% improvement threshold to prevent churn.

\subsection{Evolutionary Engine (EGPA)}

For deeper optimization, the Evolutionary Genetic Prompt Algorithm represents prompts as \texttt{PromptGenome} objects, where each genome contains a dictionary of \emph{genes}---discrete prompt components such as role definition, constraints, and style instructions.

The evolutionary cycle proceeds as follows:

\begin{algorithm}[htbp]
\caption{EGPA: Evolutionary Genetic Prompt Algorithm}
\label{alg:egpa}
\begin{algorithmic}[1]
\STATE \textbf{Input:} seed genes $g_0$, population size $P$, generations $G$
\STATE \textbf{Output:} best-performing prompt genome
\STATE $\mathcal{P} \leftarrow$ \textsc{CreatePopulation}($g_0$, $P$) \COMMENT{LLM paraphrase}
\FOR{$i = 1$ to $G$}
    \STATE \textsc{EvaluatePopulation}($\mathcal{P}$, \textit{fitness\_fn})
    \STATE $\mathcal{P}_{elite} \leftarrow$ top-$E$ from $\mathcal{P}$  \COMMENT{$E=2$}
    \STATE $\mathcal{P}' \leftarrow \mathcal{P}_{elite}$
    \WHILE{$|\mathcal{P}'| < P$}
        \STATE $p_1, p_2 \leftarrow$ \textsc{TournamentSelect}($\mathcal{P}$, $k{=}3$)
        \IF{$\text{rand}() < r_c$}  \COMMENT{$r_c = 0.7$}
            \STATE $c \leftarrow$ \textsc{UniformCrossover}($p_1$, $p_2$)
        \ELSE
            \STATE $c \leftarrow$ \textsc{Copy}($p_1$)
        \ENDIF
        \STATE \textsc{LLMMutate}($c$, $r_m{=}0.1$)  \COMMENT{Rewrite gene}
        \STATE $\mathcal{P}' \leftarrow \mathcal{P}' \cup \{c\}$
    \ENDWHILE
    \STATE $\mathcal{P} \leftarrow \mathcal{P}'$
    \IF{converged for 3 generations}
        \STATE \textbf{break}
    \ENDIF
\ENDFOR
\RETURN $\arg\max_{g \in \mathcal{P}} \text{fitness}(g)$
\end{algorithmic}
\end{algorithm}

A distinctive feature of EGPA is \emph{LLM-driven mutation}: instead of random perturbation, the LLM is prompted to ``rewrite the following text to be clearer and more effective'', producing semantically meaningful variations. Initial population diversity is achieved through LLM paraphrasing of the seed genes.

\section{Hybrid Memory System}
\label{sec:memory}

\subsection{Architecture}

The memory subsystem comprises four layers:

\textbf{Vault.} An Obsidian-style markdown vault organizes notes into subdirectories (\texttt{checkpoints/}, \texttt{derivations/}, \texttt{concepts/}). Each note includes YAML frontmatter with metadata: unique ID, type, creation timestamp, tags, wiki-style links, and relation triples.

\textbf{Knowledge Graph.} Relations are stored as subject--predicate--object triples extracted from note frontmatter. The graph supports pattern queries by any combination of subject, relation, and object, and generates derivation trees showing the provenance chain for any concept.

\textbf{Embedding Store.} ChromaDB provides persistent vector storage with cosine similarity search. When a local embedding model is available (default: Qwen3-Embedding-0.6B~\cite{qwen2024}), documents are embedded locally; otherwise, ChromaDB's default embedding function is used.

\textbf{MCP Server.} An MCP~\cite{anthropic2024mcp} server exposes five tools to the LLM: \texttt{query\_graph}, \texttt{get\_derivation\_tree}, \texttt{write\_note}, \texttt{add\_relation}, and \texttt{search\_vault}. This decoupled architecture allows the memory subsystem to be replaced without modifying the reasoning engine.

\subsection{Hybrid Retrieval with RRF}

The \texttt{Retriever} class performs hybrid search by issuing concurrent queries to the embedding store and knowledge graph, then fusing results using Reciprocal Rank Fusion~\cite{cormack2009reciprocal}:
\begin{equation}
\text{RRF}(d) = \sum_{r \in R} \frac{w_r}{k + \text{rank}_r(d)}
\label{eq:rrf}
\end{equation}
where $k{=}60$ is a smoothing constant, $R = \{\text{embedding}, \text{graph}\}$ are the retrieval systems, and $w_r$ are configurable weights ($w_{\text{emb}} {=} 0.6$, $w_{\text{kw}} {=} 0.4$ by default). Keyword extraction uses a stop-word-filtered tokenizer, limited to the top 5 terms per query.

\subsection{Cross-Session Memory}

Two additional memory mechanisms operate at the session level:

\textbf{Global Memory Store.} A persistent key-value store with category indexing enables knowledge to survive across sessions. Retrieval uses keyword-based scoring: exact key match (+10), partial key match (+5), and per-word content match (+1).

\textbf{Scratchpad.} An ordered, priority-aware working memory for the current session. Entries are formatted as markdown and injected into the system prompt. Capacity-limited (default: 100 entries) with LRU eviction.

\section{Evaluation}
\label{sec:eval}

We evaluate Cognitive Harness through a comprehensive 12-test end-to-end verification suite that covers every subsystem.

\subsection{Test Suite Design}

The test suite is organized as follows:

\begin{table}[htbp]
\caption{End-to-End Verification Test Suite}
\label{tab:tests}
\centering
\footnotesize
\begin{tabular}{@{}clp{3.8cm}@{}}
\toprule
\textbf{\#} & \textbf{Test Name} & \textbf{Coverage} \\
\midrule
1 & Vault Operations & Note CRUD, frontmatter parsing, wiki-links, relation extraction \\
2 & Knowledge Graph & Triple insertion, subject/relation queries, derivation trees \\
3 & Embedding Store & ChromaDB add/search/delete, metadata filtering \\
4 & MCP Server & Server creation, tool registration \\
5 & Hybrid Retriever & Concurrent embedding + graph search, RRF fusion \\
6 & Memory Writer & Checkpoint and relation persistence via MCP \\
7 & Configuration & YAML loading, env var resolution, default merging \\
8 & Stuck Detector & Hedging detection, repetition analysis, false-negative verification \\
9 & Plugin Registry & Dynamic registration, decorator-based discovery \\
10 & Trace Logger & Structured JSONL output with custom fields \\
11 & Orchestrator Graph & LangGraph compilation, node/edge verification \\
12 & Full Harness (LLM) & End-to-end: math, reasoning, and code understanding with real LLM \\
\bottomrule
\end{tabular}
\end{table}

\subsection{Results}

All 12 tests pass successfully. Key findings:

\textbf{Stuck Detection Accuracy.} The heuristic detector correctly identifies hedging patterns (2+ matches from 7 patterns) and 4-gram repetition ($\geq 3$ occurrences) while avoiding false positives on normal responses (Test~8). The zero-cost design (no additional LLM calls) makes detection feasible on every reasoning step.

\textbf{Hybrid Retrieval.} The RRF fusion successfully combines results from both the embedding store and the knowledge graph (Test~5), with mock evaluations confirming that both ``checkpoint''/``concept'' types and ``graph\_triple'' types appear in the fused result set.

\textbf{Graph Operations.} The knowledge graph correctly handles multi-hop derivation tree traversal with depth limiting (Test~2), and the vault's frontmatter parser correctly round-trips relation triples through YAML serialization (Test~1).

\textbf{Orchestrator Topology.} The compiled LangGraph includes all 5 core nodes (\texttt{thinker}, \texttt{stuck\_detector}, \texttt{sub\_agent\_spawner}, \texttt{compactor}, \texttt{context\_manager}) with correct conditional edges (Test~11).

\textbf{End-to-End Reasoning.} Full harness invocation with a real LLM (DeepSeek-V4-Flash) correctly solves arithmetic (``15+27=42''), multi-step word problems (``10-3+7=14''), and code understanding tasks (Test~12).

\subsection{Code Metrics}

\begin{table}[htbp]
\caption{Codebase Statistics}
\label{tab:codemetrics}
\centering
\footnotesize
\begin{tabular}{@{}lr@{}}
\toprule
\textbf{Metric} & \textbf{Value} \\
\midrule
Total Python modules & 40+ \\
Core source lines (src/harness/) & $\sim$3,200 \\
MCP server lines (mcp\_server/) & $\sim$450 \\
Test lines (tests/) & $\sim$1,800 \\
Plugin nodes & 7 (5 core + 2 memory) \\
Built-in tool categories & 6 \\
Configuration parameters & 35 \\
Dependencies & 11 \\
\bottomrule
\end{tabular}
\end{table}

\section{Discussion}
\label{sec:discussion}

\textbf{Design Trade-offs.} The heuristic stuck detector trades recall for precision---subtle reasoning loops that avoid hedging language or exact repetition will not be detected. An LLM-based detector could improve recall at the cost of additional API calls per reasoning step. The current design prioritizes latency and cost.

\textbf{Compaction Fidelity.} Two-tier compaction necessarily introduces information loss. Critical details in the working buffer may be dropped during summarization. The structured Premise$\rightarrow$Step$\rightarrow$Conclusion format mitigates this by anchoring the LLM's attention on logical structure, but formal verification of compaction correctness remains an open problem.

\textbf{Evolutionary Convergence.} EGPA's convergence depends on the quality of the fitness function and the diversity of the initial population. With the current composite score (Eq.~\ref{eq:fitness}), optimization may converge prematurely if success rate dominates. Incorporating user feedback scores (\texttt{feedback\_score} field in \texttt{PromptMetric}) could provide stronger gradient signals.

\textbf{Memory Scalability.} The global memory store uses linear-time keyword matching, which scales as $O(nm)$ where $n$ is the number of entries and $m$ is the query length. For vaults exceeding $\sim$10,000 entries, the embedding-based retrieval path should be preferred. The configurable \texttt{max\_entries} parameter (default 10,000) provides a practical bound.

\textbf{MCP Interoperability.} The use of MCP for memory access provides protocol-level interoperability---the vault server can be replaced with any MCP-compatible knowledge base. However, the current MCP client management requires subprocess lifecycle handling, which adds complexity on Windows platforms.

\section{Conclusion}
\label{sec:conclusion}

We have presented Cognitive Harness, a plugin-based LLM reasoning engine that addresses the fundamental problem of reasoning stalls in multi-step LLM interactions. Through the integration of metacognitive stuck detection, parallel sub-agent recovery, two-tier context compaction, and evolutionary prompt optimization, the system maintains productive reasoning over extended sessions. The hybrid memory subsystem with RRF-fused retrieval enables knowledge persistence and cross-session learning, while the MCP-based architecture ensures modularity and interoperability.

The system's plugin architecture---where every component from the thinker to the stuck detector is a self-registering \texttt{BaseNode}---enables targeted replacement of individual components. This design facilitates future work including: (1)~replacing the heuristic stuck detector with a learned classifier trained on labeled stall examples; (2)~incorporating semantic similarity alongside n-gram repetition for more nuanced loop detection; (3)~extending EGPA with multi-objective Pareto optimization for the fitness function; and (4)~scaling the knowledge graph to support distributed, multi-vault federations.

Cognitive Harness demonstrates that the combination of metacognitive monitoring, parallel recovery, context management, and self-improvement can transform LLMs from single-turn response generators into sustained reasoning engines. The complete source code is publicly available as open-source software.

\begin{thebibliography}{00}

\bibitem{achiam2023gpt4}
J.~Achiam \textit{et al.}, ``GPT-4 technical report,'' \textit{arXiv preprint arXiv:2303.08774}, 2023.

\bibitem{deepseekai2024deepseekv3}
DeepSeek-AI, ``DeepSeek-V3 technical report,'' \textit{arXiv preprint arXiv:2412.19437}, 2024.

\bibitem{anthropic2024claude}
Anthropic, ``The Claude model family,'' Anthropic, Technical Report, 2024.

\bibitem{wei2022chain}
J.~Wei \textit{et al.}, ``Chain-of-thought prompting elicits reasoning in large language models,'' in \textit{Proc. NeurIPS}, 2022, pp.~24824--24837.

\bibitem{langchain2023}
H.~Chase, ``LangChain: Building applications with LLMs through composability,'' 2023. [Online]. Available: https://github.com/langchain-ai/langchain

\bibitem{langgraph2024}
LangChain, ``LangGraph: Build language agents as graphs,'' 2024. [Online]. Available: https://github.com/langchain-ai/langgraph

\bibitem{khattab2023dspy}
O.~Khattab \textit{et al.}, ``DSPy: Compiling declarative language model calls into self-improving pipelines,'' in \textit{Proc. ICLR}, 2024.

\bibitem{yao2023react}
S.~Yao \textit{et al.}, ``ReAct: Synergizing reasoning and acting in language models,'' in \textit{Proc. ICLR}, 2023.

\bibitem{shinn2023reflexion}
N.~Shinn \textit{et al.}, ``Reflexion: Language agents with verbal reinforcement learning,'' in \textit{Proc. NeurIPS}, 2023.

\bibitem{wu2023autogen}
Q.~Wu \textit{et al.}, ``AutoGen: Enabling next-gen LLM applications via multi-agent conversation,'' \textit{arXiv preprint arXiv:2308.08155}, 2023.

\bibitem{crewai2024}
J.~Moura, ``CrewAI: Framework for orchestrating role-playing autonomous AI agents,'' 2024. [Online]. Available: https://github.com/crewAIInc/crewAI

\bibitem{fernando2023promptbreeder}
C.~Fernando \textit{et al.}, ``PromptBreeder: Self-referential self-improvement via prompt evolution,'' \textit{arXiv preprint arXiv:2309.16797}, 2023.

\bibitem{zhou2023large}
Y.~Zhou \textit{et al.}, ``Large language models are human-level prompt engineers,'' in \textit{Proc. ICLR}, 2023.

\bibitem{yao2024tree}
S.~Yao \textit{et al.}, ``Tree of thoughts: Deliberate problem solving with large language models,'' in \textit{Proc. NeurIPS}, 2024.

\bibitem{madaan2023selfrefine}
A.~Madaan \textit{et al.}, ``Self-Refine: Iterative refinement with self-feedback,'' in \textit{Proc. NeurIPS}, 2023.

\bibitem{packer2023memgpt}
C.~Packer \textit{et al.}, ``MemGPT: Towards LLMs as operating systems,'' \textit{arXiv preprint arXiv:2310.08560}, 2023.

\bibitem{lewis2020rag}
P.~Lewis \textit{et al.}, ``Retrieval-augmented generation for knowledge-intensive NLP tasks,'' in \textit{Proc. NeurIPS}, 2020, pp.~9459--9474.

\bibitem{cormack2009reciprocal}
G.~V.~Cormack, C.~L.~A.~Clarke, and S.~B{\"u}ttcher, ``Reciprocal rank fusion outperforms Condorcet and individual rank learning methods,'' in \textit{Proc. SIGIR}, 2009, pp.~758--759.

\bibitem{anthropic2024mcp}
Anthropic, ``Model Context Protocol specification,'' 2024. [Online]. Available: https://modelcontextprotocol.io

\bibitem{qwen2024}
Qwen Team, ``Qwen3 technical report,'' Alibaba Group, 2024.

\end{thebibliography}

\end{document}
